\section{Fase 4. Accionamiento electrónico, operación manual y referenciación}
\label{sec:discusion_fase4}

La implementación progresiva separó el desarrollo del accionamiento, la lectura de los límites y la integración del mando. El movimiento por comandos proporcionó una base para incorporar posteriormente la interacción inalámbrica. Esta secuencia correspondió al orden de desarrollo descrito por el autor y no al orden de arranque del firmware integrado, que también incorporó funciones de etapas posteriores.

La distribución entre placas se describió según sus responsabilidades. La ESP32 concentró la interacción con el usuario, la OLED y los servomotores; la Portenta concentró el accionamiento de los ejes, los límites, la calibración y los estados generales. Aunque las órdenes del operador se originaron en la ESP32, las transacciones I2C fueron iniciadas por la Portenta. Por ello, la procedencia de una orden de movimiento no equivalió al papel de maestro del bus ni a la coordinación general del programa.

La separación entre el bus de la pantalla y el enlace de control permitió organizar ambos intercambios de manera independiente. Además, el firmware incorporó reintentos de inicialización de la OLED sin detener deliberadamente la atención del resto de las funciones. El intercambio de estados permitió que la interfaz presentara las direcciones ejecutadas y los límites informados por la Portenta, en lugar de limitarse a mostrar la entrada del joystick. No obstante, esos estados no constituyeron una medición independiente de la posición física del efector final.

El formato fijo de los paquetes estableció una representación común de las órdenes y de los estados en ambos programas. El CRC aportó una comprobación de integridad, mientras que el número de secuencia y la sesión de arranque permitieron distinguir la falta de actualización y el reinicio de la ESP32. Estas comprobaciones atendieron condiciones diferentes: la recepción de un mensaje con CRC válido no aseguró, por sí sola, que sus datos se hubieran actualizado recientemente. Del mismo modo, la presencia de CRC no garantizó la detección de cualquier alteración posible.

La entrega de una copia previamente preparada y la validación diferida de los mensajes recibidos separaron el procesamiento de datos de las funciones inmediatas de atención I2C. La prioridad asignada al envío de estados favoreció la actualización de la interfaz dentro de la planificación implementada, aunque pudo posponer un sondeo de control hasta la siguiente iteración. Sin mediciones del enlace, estos mecanismos no demostraron una latencia determinada ni ausencia de pérdidas de comunicación.

La zona neutra del mando delimitó la región en la que no se ordenó movimiento. Fuera de ella se seleccionó una dirección discreta, por lo que la amplitud del joystick no estableció una velocidad proporcional de traslación. Esta implementación atendió el movimiento manual por eje y aportó las funciones previstas en el segundo objetivo específico, sin demostrar por sí sola una mejora cuantificada del mantenimiento o del posicionamiento.

La doble aproximación a los finales estableció una secuencia definida de referenciación. El conteo entre extremos y la posterior ubicación en el centro proporcionaron un origen común para expresar los desplazamientos. Sin embargo, el conteo de pulsos generados representó la posición calculada por el programa y no una medición continua independiente del desplazamiento real del brazo.

La función denominada cinemática inversa correspondió a una conversión cartesiana de X/Y a destinos de pasos mediante factores de escala. Su alcance quedó limitado por los ejes efectivamente accionados: la recepción de un argumento Z no implicó posicionamiento cartesiano tridimensional. Asimismo, el movimiento por joystick no utilizó esa conversión, por lo que no se presentó como una ejecución de cinemática inversa mediante el mando.

La confirmación del recorrido físico de Y en 336~mm estableció la correspondencia entre el dato mecánico y la longitud utilizada para calcular el factor de pasos por milímetro. Esta correspondencia no constituyó por sí sola una validación de exactitud: la posición calculada también dependió del conteo entre extremos y de que el desplazamiento físico correspondiera a los pulsos generados. No se atribuyeron valores de precisión o repetibilidad a la rutina sin registros de posiciones obtenidas en ensayos.

Las comprobaciones de límites, la neutralización de órdenes y la validación de paquetes incorporaron respuestas de software frente a determinadas condiciones de fallo. Su presencia en el código no demostró una certificación de seguridad ni confirmó el comportamiento físico de todas las conexiones. La validación de las protecciones y del posicionamiento requirió distinguir los resultados observados en el prototipo de las funciones previstas en el firmware.
